\documentclass[11pt,a4paper,landscape]{article}
\usepackage[utf8]{inputenc}
\usepackage[T1]{fontenc}
\usepackage[czech]{babel}
\usepackage{lmodern}
\usepackage[margin=1.2cm]{geometry}
\usepackage{tikz}
\usetikzlibrary{shapes.geometric,positioning}
\pagestyle{empty}

\begin{document}
\begin{center}
\begin{tikzpicture}[
  font=\small,
  center/.style={
    ellipse,
    draw=blue!60!black,
    fill=blue!12,
    very thick,
    minimum width=6cm,
    minimum height=2.6cm,
    align=center,
    font=\Large\bfseries
  },
  branch/.style={
    rectangle,
    rounded corners=4pt,
    draw=blue!55!black,
    fill=blue!8,
    thick,
    minimum width=3.8cm,
    minimum height=1.2cm,
    align=center,
    font=\bfseries
  },
  detail/.style={
    rectangle,
    rounded corners=2pt,
    draw=gray!55,
    fill=white,
    minimum width=3.4cm,
    minimum height=0.9cm,
    align=center,
    font=\footnotesize
  },
  edge/.style={thick, draw=blue!45!black}
]

% Pozice středu, větví a detailů
\coordinate (c) at (0,0);

\coordinate (b1) at (-5.6,2.6);
\coordinate (b2) at (-5.6,-2.6);
\coordinate (b3) at (5.6,2.6);
\coordinate (b4) at (5.6,-2.6);

\coordinate (d11) at (-9.8,4.1);
\coordinate (d12) at (-9.8,1.1);
\coordinate (d21) at (-9.8,-1.1);
\coordinate (d22) at (-9.8,-4.1);

\coordinate (d31) at (9.8,4.1);
\coordinate (d32) at (9.8,1.1);
\coordinate (d41) at (9.8,-1.1);
\coordinate (d42) at (9.8,-4.1);

% Spojení se kreslí před uzly, aby čáry nebyly vidět přes výplně boxů
\draw[edge] (c) -- (b1);
\draw[edge] (c) -- (b2);
\draw[edge] (c) -- (b3);
\draw[edge] (c) -- (b4);

\draw[edge] (b1) -- (d11);
\draw[edge] (b1) -- (d12);
\draw[edge] (b2) -- (d21);
\draw[edge] (b2) -- (d22);

\draw[edge] (b3) -- (d31);
\draw[edge] (b3) -- (d32);
\draw[edge] (b4) -- (d41);
\draw[edge] (b4) -- (d42);

% Uzly se kreslí až po spojích
\node[center] at (c) {TÉMA\\OTÁZKY};

\node[branch] at (b1) {ČÁST 1};
\node[branch] at (b2) {ČÁST 2};
\node[branch] at (b3) {ČÁST 3};
\node[branch] at (b4) {ČÁST 4};

\node[detail] at (d11) {detail A};
\node[detail] at (d12) {detail B};
\node[detail] at (d21) {detail C};
\node[detail] at (d22) {detail D};

\node[detail] at (d31) {detail E};
\node[detail] at (d32) {detail F};
\node[detail] at (d41) {detail G};
\node[detail] at (d42) {detail H};

\end{tikzpicture}
\end{center}
\end{document}
